\begin{outline}[enumerate]

\1 \textbf{What Is Logic?}
\\ Required Reading: \cite[Chapters 1-5 (pp. 2-48)]{Mag25} \\
Digging Deeper: \cite[Chapters 1-4,6-7,11 (pp. 1-39,44-58,94-104)]{Smi20} \\
Focus:
informal presentation of \emph{validity}, \emph{soundness}, \emph{entailment},
\emph{satisfiability}; distinguishing \emph{deductive} from \emph{inductive}
arguments; translating ordinary English into truth-functional logic (TFL);
understanding the meaning of TFL connectives;
\1 \textbf{Truth and Grammar in TFL}
\\ Required Reading: \cite[Chapters 6-11 (pp. 49-84)]{Mag25} \\
Digging Deeper: \cite[Chapters 8-13,18-19 (pp. 59-120,148-172)]{Smi20} \\
Focus:
\emph{sentences} in TFL, \emph{parse trees}, connective \emph{scope}, types of
\emph{ambiguity}; understanding \emph{truth-functional} connectives; building
\emph{truth tables} for complex sentences;
\1 \textbf{Semantic Strategies in TFL}
\\ Required Reading: \cite[Chapters 12-15 (pp. 85-108)]{Mag25} \\
Digging Deeper: \cite[Chapters 14-16 (pp. 120-137)]{Smi20} \\
Focus:
testing logical properties with truth tables; efficiency strategies when using
truth tables; using partial tables for validation;\\
\1 \textbf{Natural Deduction in TFL}
\\ Required Reading: \cite[Chapter 16-17 (pp. 110-141)]{Mag25} \\
Digging Deeper: \cite[Chapters 20-22 (pp. 172-211)]{Smi20} \\
Focus:
presentation of the Fitch-style \emph{proof} system; \emph{introduction} and
\emph{elimination} rules for TFL connectives;
\1 \textbf{Proofs in TFL}
\\ Required Reading: \cite[Chapter 18 (pp. 142-161)]{Mag25} \\
Digging Deeper: \cite[Chapter 17 (pp. 145-147)]{Smi20} \\
Focus:
proof strategies; working backward from the conclusion; working forward from
the premises; building and exploiting \emph{subproofs};
\1 \textbf{Additional Rules}
\\ Required Reading: \cite[Chapters 19-21 (pp. 162-181)]{Mag25} \\
Digging Deeper: \cite[Chapter 23 (pp. 211-215)]{Smi20} \\
Focus:
extending the basic system with additional rules;
\1 \textbf{Basic TFL Metalogic}
\\ Required Reading: \cite[Chapter 22 (pp. 182-188)]{Mag25} \\
Digging Deeper: \cite[Chapter 24 (pp. 216-228)]{Smi20} \\
Focus:
connecting \emph{proof theory} to the semantics of \emph{truth-functional}
logic;
\1 \textbf{An Informal Survery of First-Order Logic}
\\ Required Reading: \cite[Chapter 23-27 (pp. 190-242)]{Mag25} \\
Digging Deeper: \cite[Chapters 25-27,31 (pp. 228-258,293-298)]{Smi20} \\
Focus:
translating simple \emph{universal} and \emph{existential} sentences;
\emph{predicate} notation; nested \emph{quantifiers}; the \emph{identity}
predicate; \emph{definite descriptions};
\1 \textbf{Grammar of FOL}
\\ Required Reading: \cite[Chapter 28-29 (pp. 243-256)]{Mag25} \\
Digging Deeper: \cite[Chapters 28-30 (pp. 258-289)]{Smi20} \\
Focus:
\emph{sentences} and \emph{scope ambiguities} in first-order logic;
\1 \textbf{FOL Semantics}
\\ Required Reading: \cite[Chapter 30-35 (pp. 258-295)]{Mag25} \\
Focus:
assessing the truth of FOL sentences relative to an \emph{interpretation};
testing validity via \emph{counter-models}
\1 \textbf{FOL Proofs}
\\ Required Reading: \cite[Chapters 36-41 (pp. 297-335)]{Mag25} \\
Digging Deeper: \cite[Chapters 32-42, Appendix (pp. 299-412)]{Smi20} \\
Focus:
introduction and elimination rules for universal and existential quantifiers;
identity rules; connecting FOL proof theory to interpretations;

\end{outline}
